\answer{ex:21:1}{(а)~$T=\tau_{\text{відн}}\ln\frac{1}{1-x^*}$: $1.84$~с і $3.68$~с. (б)~Від $3.36$ до $4.24$~с: чутливість $\dd T/\dd x^*=\tau_{\text{відн}}/(1-x^*)=80$~с, при $x^*\to1$ залпи нерегулярні.}
\answer{ex:21:2}{(а)~$8.6$~Вт, як у~\eqref{eq:energy}. (б)~$205$~Мбіт/с, $0.2$~Вт --- у $2\cdot10^3$ разів більше за культуру ($10^{-4}$~Вт).}
\answer{ex:21:3}{$x_{\text{уст}}=1/(1+Ur_b\tau_{\text{відн}})<x^*$ при $r_b>(1/x^*-1)/(U\tau_{\text{відн}})=0.28$~Гц для $x^*=0.9$ і $0.025$~Гц для $x^*=0.99$; сила синапсу падає до $x_{\text{уст}}$, тобто на $10\%$ і $1\%$.}
\answer{ex:21:4}{$u(t-k)$ ортонормовані, $m_k=\|P_Vu(t-k)\|^2$, де $P_V$ --- проєктор на $N$-вимірну оболонку виходів; за нерівністю Бесселя $\sum_k m_k\le N$.}
\answer{ex:21:5}{$\mathrm{MC}\approx9$--$11$ з $\tanh$ і $15$--$18$ для лінійного резервуара (три затравки), обидва $\le30$: нелінійність платить пам'яттю.}
\answer{ex:21:6}{(а)~$n\le102$. (б)~На $5.6$.}
\answer{ex:21:7}{Відкрита задача. Ключовий контроль: заморозити зчитування і перевірити, чи зберігається покращення після паузи без стимуляції; якщо ні, змінився стан, а не ваги.}
